% cap03/3.5_variables_indicadores.tex
\section{Variables e Indicadores}

Para despojar a este trabajo ingenieril de cualquier sesgo subjetivo y alcanzar la pureza del an\'alisis algor\'itmico, las variables de evaluaci\'on se anclan de forma dogm\'atica a m\'etricas estandarizadas de calidad de software y pruebas automatizadas. No se admiten apreciaciones cualitativas; cada indicador se traduce a una f\'ormula algor\'itmica expl\'icita.

\subsection{Matriz de Operacionalizaci\'on: Funcionamiento y Usabilidad}

\textbf{Variable Independiente: La Arquitectura \textit{Democra}.} La arquitectura de \textit{Democra} (compuesta por \textit{React.js}, \textit{Node.js}, y \textit{PostgreSQL} con RLS) se erige como la causa motriz. Su operaci\'on es inyectada sobre el entorno de pruebas de caja blanca.

\textbf{Variables Dependientes: Funcionamiento y Usabilidad.} En coherencia con el
resumen del estudio, la evaluaci\'on se estructura en torno a dos variables
dependientes, cada una operacionalizada mediante indicadores objetivos:
\begin{itemize}
    \item \textbf{Funcionamiento:} capacidad del sistema para operar de forma
    correcta, eficiente y segura. Se operacionaliza mediante la eficiencia de
    desempe\~no (comportamiento temporal bajo carga), el bloqueo de fallos de
    concurrencia (efectividad del aislamiento por RLS) y la cobertura de c\'odigo
    alcanzada por las pruebas automatizadas.
    \item \textbf{Usabilidad:} grado de facilidad con que los operadores emplean el
    sistema. Se operacionaliza mediante la escala estandarizada \textit{System
    Usability Scale} (SUS), aplicada a la muestra de operadores y reportada en el
    cap\'itulo de resultados.
\end{itemize}

\subsection{Traducci\'on Matem\'atica y Umbrales R\'igidos}

Para la \textbf{Eficiencia de Desempe\~no}, se establece la m\'etrica de \textit{Latencia Media de Transacci\'on As\'incrona} ($L_m$). La f\'ormula algor\'itmica es la sumatoria de las diferencias de tiempo ($\Delta t$) entre el env\'io del \textit{WebSocket} y la recepci\'on del \textit{ACK}, dividida entre el n\'umero total de transacciones ($N$).
\begin{equation}
L_m = \frac{1}{N} \sum_{i=1}^{N} (t_{ACK\_i} - t_{Send\_i})
\end{equation}
Se impone un \textbf{umbral de aceptaci\'on r\'igido}: $L_m < 200 \text{ ms}$. Cualquier resultado superior se considera un fallo arquitect\'onico inaceptable.

Asimismo, se define la m\'etrica de \textbf{Cobertura de Bloqueo de Fallos de Concurrencia} ($C_b$).
\begin{equation}
C_b = \left( \frac{\text{Ataques Bloqueados por RLS}}{\text{Total de Ataques Inyectados}} \right) \times 100\%
\end{equation}
El umbral absoluto es del $\mathbf{100\%}$. No se tolera la m\'as m\'inima fuga de datos.

Para la \textbf{Cobertura de C\'odigo}, la m\'etrica de porcentaje global de cobertura ($Cov$) se expresa como la proporci\'on de l\'ineas de c\'odigo ($L_c$) transitadas exitosamente por las suites de prueba \textit{Jest} respecto al total de l\'ineas ejecutables del proyecto ($L_t$).
\begin{equation}
Cov = \left( \frac{L_c}{L_t} \right) \times 100\%
\end{equation}
El umbral te\'orico de aceptaci\'on es una cobertura global $\ge 85\%$, considerado el est\'andar de alta calidad en la industria del software para despliegues a producci\'on.
